\section{Discussion, Limitations and Future Work}
\label{sec:discussion}

\subsection{Interpretation}

Three findings carry beyond the specific model.

First, \textbf{graph structure is not automatically informative}, even when the
data-generating process contains explicit spillover. The static intercept
collapse of Section~\ref{subsec:collapse} shows that spillover which is
permanent and early is absorbed into node-level state and becomes
observationally equivalent to a fixed effect. A benchmark constructed that way
will report that graph models do not help, and the report will be correct about
the benchmark while saying nothing about the method. Practitioners evaluating
spatio-temporal architectures should verify that their spillover process
remains dynamic \emph{within the evaluation window}, and we would encourage
treating the shuffled-edge ablation as a standard diagnostic rather than an
optional extra: it was the shuffled condition, not the edgeless one, that most
clearly separated the working benchmark from the degenerate one.

Second, \textbf{the advantage of nonlinear message passing is
regime-dependent}. In the untreated equilibrium a linear spatial lag is
statistically indistinguishable from \CGCEN{} ($p = 0.127$). Under active
intervention the gap opens to roughly $3.3$ MAE with large effect sizes. The
practical reading is that a planning agency running a steady-state monitoring
model has little reason to prefer a graph neural network over classical spatial
econometrics~\cite{anselin1988,lesage2009}, whereas one comparing prospective
interventions with spatially propagating consequences does.

Third, \textbf{index design can silently determine conclusions}. The
growth-rate Displacement Pressure Index is internally coherent and, applied to
a permanent policy shift, reports approximately zero effect. An analyst who did
not interrogate the functional form could reasonably conclude that inclusionary
housing has no displacement consequence, when in fact the index is structurally
incapable of registering a sustained level difference. The level-aware
correction of \eqref{eq:burden} changes the measured contrast by two orders of
magnitude without altering a single simulated quantity.

\subsection{Limitations}

\subsubsection{No full-scale empirical evaluation}

This is the principal limitation and we state it plainly: \textbf{all
quantitative results in this paper are derived from synthetic data.} We
implement a complete empirical pipeline---a loader that ingests a census-tract
table with ACS and PLUTO fields (median gross rent B25064, median household
income B19013, population B01003, housing units B25001, assessed land value,
lot and residential area, and a derived transit accessibility score), builds
tract adjacency from centroid geodesics by $k$-nearest neighbours under a
haversine metric, and emits a PyTorch Geometric graph---and we validate it end
to end. But we validate it on a \emph{schema-compatible sample}, not on a real
extract. The loader reports its provenance through an explicit
\texttt{is\_synthetic\_sample} flag precisely so that this distinction cannot
be lost downstream, and we have not presented sample-derived numbers anywhere
in this manuscript.

Consequently the results establish \emph{internal} validity---given a
data-generating process with known dynamic spillover, the architecture recovers
it, and we can say exactly when it fails to. They establish nothing about
external validity. Real tract panels are irregular, sparsely observed, subject
to boundary changes across decennial revisions, and contain confounding that
our generator does not simulate. Whether the conditional advantage we document
survives contact with such data is untested, and we would regard any policy
application of \CGCEN{} prior to that test as premature.

\subsubsection{Scale and structure of the synthetic city}

The $4 \times 4$ lattice with $16$ units and $18$ steps yields nine training
windows. This is small, and it interacts with the results in a specific way:
model capacity is cheap to waste, which is part of why the Temporal MLP is so
variable ($\pm 7.97$ MAE) and why the linear SLM is competitive at all.
Real metropolitan graphs are larger, irregular, and have heavy-tailed degree
distributions that a regular lattice does not exercise. We do not claim the
observed ordering transfers unchanged to that setting.

\subsubsection{Simplified interventions and known ground truth}

Policies are applied as multiplicative adjustments on fixed node sets with
fixed timing. Real interventions are heterogeneous in intensity, phased,
politically endogenous to the neighbourhoods that receive them, and often
anticipated by markets before enactment. Treatment endogeneity in particular
would break the exogeneity our design assumes, and addressing it requires
identification machinery---instruments, difference-in-differences with spatial
lags, or an explicit assignment model---that we have not incorporated.

\subsubsection{Partial interference and the exposure summary}

We assume interference operates only through the degree-normalised
one-hop exposure of \eqref{eq:exposure}. This rules out longer-range effects
that are not mediated by intervening units, and it compresses the neighbourhood
assignment vector into a scalar. Both are standard simplifications in the
interference literature~\cite{hudgens2008,forastiere2021}, and both are
falsifiable: richer exposure mappings would be a direct extension.

\subsection{Future Work}

The immediate priority is \textbf{empirical replication}. Sourcing a genuine
NYC ACS/PLUTO tract extract and re-running the benchmark against it is the
single highest-value next step, and the loader exists to make it a
configuration change rather than an engineering project. A useful intermediate
would be a semi-synthetic design: real tract geometry, adjacency and initial
conditions, with simulated interventions layered on, which preserves ground
truth while introducing realistic spatial structure.

Beyond that, three directions follow naturally. \emph{Richer exposure
mappings} would relax the one-hop scalar assumption toward multi-hop or
distance-weighted summaries. \emph{Uncertainty quantification} over the
estimated ITE and STE would let the scorecard report credible intervals on
policy contrasts rather than point estimates; conformal or ensemble approaches
are the obvious candidates. \emph{Treatment assignment modelling} would move
the framework from the ex-ante comparative setting, where assignment is
stipulated, toward observational evaluation, where it must be estimated.

\subsection{Conclusion}

We presented CivicTwin, a counterfactual graph learning framework that
decomposes urban policy effects into a direct term and a spatially
interfering term under partial interference, with the split enforced
architecturally rather than through regularisation. On a controlled synthetic
benchmark across ten seeds and four policy scenarios, \CGCEN{} outperforms a
spatial lag model, a graph-blind temporal MLP and naive persistence, and
topology ablations confirm the advantage is carried by the graph: randomly
rewiring the adjacency is more damaging than removing it entirely.

The advantage is conditional, and we have tried to be precise about the
condition. \CGCEN{} separates from classical spatial econometrics under active
policy intervention ($p = 0.0015$--$0.0033$) but not in the untreated
equilibrium ($p = 0.127$). Under an earlier generator in which spillover
arrived as a single permanent shock, it separated from nothing at all---no
model benefited from the graph, because permanent early spillover collapses
into a static intercept. We report that negative result alongside the positive
one because the boundary between the two regimes is, for anyone building
spatio-temporal policy models, more actionable than either result alone.
